\section{What would justify handing over a release?}\label{sec:intro}
An engineering team gives a software agent a build environment. The agent can reconcile an inventory, prepare a patch and run the test suite. A production release requires a separate authorization. The team would like to entrust real work to the agent while preserving that requirement. Repeating staging tests is easy. Establishing who controls the final production action is harder.

Imagine two installations. Both display an approval prompt and reject every unsigned staging request. In one, the production finalizer checks the signed grant on every route. In the other, one admitted route goes through an older finalizer. Both explanations fit the observations. A verified read of the production route and its guard could settle the deployment question while leaving the compiler, the agent's reasoning and much of the installation unknown. The inventory and test work could also be entrusted under a common restriction while that question is open. This is a running illustration. The measured services appear in Section~\ref{sec:experiments}.

Three decisions are involved. \emph{Specification} fixes what counts as an unauthorized release, which routes are in scope and what useful work is owed. \emph{Identification} asks which answers remain possible given the evidence. \emph{Realization} asks whether a controller can meet the specification across the possibilities that remain. A separate working plan for each installation does little for a controller that must choose before learning which installation it operates.

These distinctions matter for increasingly capable systems. A test can change the incentives of its subject. An approval can change what the subject is allowed to do. A dashboard can miss the very operation being counted. Even an always-safe policy can be useless if it declines every job. Each difficulty changes what evidence would warrant a deployment decision. We ask how to find that evidence, what it costs under a given interface and how a controller can use it without access to hidden state.

The contribution is a chain of checkable claims. The original Perceptive Closure manuscript and subsequent theorem charter frame the construction \citep{d35,charter}. It measures which source-level resources a repair needs. We lift this construction to a set of possible systems, compute the minimum reads that settle a property or make a policy feasible, and state what connects evaluated effects to actual deployment effects. One family separates the price of knowing a target from that of recovering an installation. A second gives a common policy that pays for observations, withstands a specified response to its verdict and performs real work. The proofs are conditional on the source description, query rights, outcomes and costs. The experimental record includes settings where ordinary full acquisition, direct observation or a simple matched solver did as well or better.

\figpc{worlds}{0.91}{A rehearsal admits four possible installations in the running example. A genuine read of the operative guard retains two. They use different build backends and agree on the release guarantee. The example assumes independently established coverage of every admitted production route.}{fig:worlds}
\noindent A practical question runs through the paper. Which difference between two installations could change the answer, and which available observation or lawful action would expose that difference?
